\documentclass[11pt,reqno]{amsart}

\usepackage[T1]{fontenc}
\usepackage{lmodern}
\usepackage{microtype}
\usepackage{mathtools}
\usepackage{amssymb}
\usepackage{booktabs}
\usepackage{array}
\usepackage{enumitem}
\usepackage{xcolor}
\usepackage{xurl}
\usepackage[colorlinks=true,linkcolor=blue!45!black,citecolor=blue!45!black,urlcolor=blue!45!black,
  pdftitle={The bordered Jacobian of binary-form multiplication: an integral foundation for the degree-difference principle},
  pdfauthor={Anonymous},
  pdfsubject={Binary-form factorisation spaces and the three-dimensional Keller counterexample},
  pdfkeywords={Jacobian conjecture, Keller map, binary forms, resultant, Sylvester matrix, Koszul complex, Vandermonde determinant, machine verification}]{hyperref}
\usepackage{aliascnt}
\usepackage[nameinlink,capitalise,noabbrev]{cleveref}

\setlength{\textwidth}{6.15in}
\setlength{\textheight}{9.0in}
\setlength{\oddsidemargin}{0.18in}
\setlength{\evensidemargin}{0.18in}
\setlength{\topmargin}{-0.15in}
\setlength{\parskip}{2pt}
\raggedbottom

\newtheorem{theorem}{Theorem}[section]

\newaliascnt{proposition}{theorem}
\newtheorem{proposition}[proposition]{Proposition}
\aliascntresetthe{proposition}

\newaliascnt{corollary}{theorem}
\newtheorem{corollary}[corollary]{Corollary}
\aliascntresetthe{corollary}

\newaliascnt{lemma}{theorem}
\newtheorem{lemma}[lemma]{Lemma}
\aliascntresetthe{lemma}

\theoremstyle{definition}
\newaliascnt{definition}{theorem}
\newtheorem{definition}[definition]{Definition}
\aliascntresetthe{definition}

\theoremstyle{remark}
\newaliascnt{remark}{theorem}
\newtheorem{remark}[remark]{Remark}
\aliascntresetthe{remark}

\crefname{theorem}{Theorem}{Theorems}
\crefname{proposition}{Proposition}{Propositions}
\crefname{corollary}{Corollary}{Corollaries}
\crefname{lemma}{Lemma}{Lemmas}
\crefname{definition}{Definition}{Definitions}
\crefname{remark}{Remark}{Remarks}

\newcommand{\A}{\mathbb A}
\newcommand{\C}{\mathbb C}
\newcommand{\Z}{\mathbb Z}
\newcommand{\Q}{\mathbb Q}
\newcommand{\OO}{\mathcal O}
\newcommand{\Cl}{\operatorname{Cl}}
\newcommand{\Res}{\operatorname{Res}}
\newcommand{\Sym}{\operatorname{Sym}}
\newcommand{\Vand}{\operatorname{V}}
\newcommand{\rank}{\operatorname{rank}}
\newcommand{\abs}[1]{\lvert #1\rvert}
\newcommand{\set}[1]{\left\{#1\right\}}

\title[The bordered Jacobian of binary-form multiplication]
{The bordered Jacobian of binary-form multiplication:\\
an integral foundation for the degree-difference principle}

% Replace ``Anonymous'' before public circulation if appropriate.
\author{Anonymous}
\date{6 August 2026 (v0.3-candidate)}
\thanks{Version 0.3-candidate.  Prepared for the Evidence Press
evidence bundle; the verification suite, receipts, environment
records, search reports and review correspondence ship with the
release archive.}

\subjclass[2020]{Primary 13P15; Secondary 14R15, 15A15, 15A54}
\keywords{Jacobian conjecture, Keller map, binary forms, resultant, Sylvester matrix, Koszul complex, Vandermonde determinant, machine verification}

\begin{document}

\begin{abstract}
Let $V_d$ denote the binary forms of degree $d$ and let
$m(A,B)=AB$ be multiplication $V_r\times V_s\to V_{r+s}$.  Writing
$M$ for the Jacobian matrix of $m$ in coefficient coordinates and
$\kappa=(a_0,\dots,a_r,-b_0,\dots,-b_s)$ for the relative-scaling
vector, we prove the \emph{bordered Jacobian identity}: for every row
vector $v$,
\[
 \det\begin{pmatrix}M\\ v\end{pmatrix}
 =(-1)^{s(r+1)+1}\,\Res(A,B)\,\langle v,\kappa\rangle
 \qquad\text{in }\Z[a,b,v],
\]
valid for all $r,s\ge1$ \emph{including} $r=s$.  The proof is
elementary: interpolation at the roots plus one auxiliary
point, a gap-Vandermonde identity, and the product formula.  It uses
no irreducibility or factoriality input, computes every sign, and,
although it argues over $\C$, it delivers an integral identity valid
after arbitrary base change.

Contracting the border against the Euler derivative of any
bihomogeneous $g$ of weight $(p,q)$ gives
$\det D(m,g)=(-1)^{s(r+1)+1}(p-q)\,g\Res$.  The choice $g=\Res$
recovers the degree-difference identity
$\det D\Phi_{r,s}=(-1)^{s(r+1)}(r-s)\Res^2$ of the July 2026
release, now over $\Z$; reduction shows $\Phi_{r,s}$ is
\'etale on the coprime locus over any field in which
$(s-r)\cdot1$ is invertible.  The choice $g=h\circ m$ gives zero: no function of the
product can complete the chart.  At $r=s$ the bordered identity does
not degenerate; the collapse of $\Phi_{r,r}$ is purely the vanishing
of the Euler weight $s-r$.  In characteristic zero the normalisers
are classified exactly: $(m,g)$ is \'etale on the whole coprime
locus if and only if $\delta g=c\Res^n$ for the scaling derivation
$\delta$, some constant $c\ne0$ and $n\ge1$---equivalently, $g$ is
a constant multiple of a power of $\Res$ plus an element of
$\ker\delta$; in characteristic $p$ the same holds when
$p\nmid n(s-r)$.

An accompanying exact suite verifies the identities of this paper per
bidegree---the central identity in two separately written
computational backends (SymPy and FLINT)---through $r+s\le8$ with spot checks to
$(5,6)$, with negative controls.  We also give a formalization
blueprint: a monic analogue of the bordered identity is already
machine-checked inside Mathlib.
\end{abstract}

\maketitle

\section{Introduction}

The July 2026 release
\emph{The degree-difference principle and affine slices of binary-form
factorisation spaces} \cite{DDAS} organised the geometry behind the
three-dimensional Keller counterexample
\cite{AlpogeAnnouncement,AnonymousCounterexample} around one
determinant identity.  For binary forms $A$ of degree $r$ and $B$ of
degree $s$, with
\[
 \Phi_{r,s}:V_r\times V_s\longrightarrow V_{r+s}\times\A^1,
 \qquad (A,B)\longmapsto\bigl(AB,\Res(A,B)\bigr),
\]
it asserted
\begin{equation}\label{eq:degree-difference}
 \det D\Phi_{r,s}=(-1)^{s(r+1)}(r-s)\Res(A,B)^2 .
\end{equation}
The proof given there ran through two fragile steps: the
irreducibility of the resultant hypersurface combined with a
unique-factorisation weight argument, to show the determinant is a
constant multiple of $\Res^2$; and a hand count of permutation
inversions at the base point $(X^r,Y^s)$, to pin the constant.  Its
machine ledger checked the base-point sign matrix and the
cubic-specific algebra, but not identity
\eqref{eq:degree-difference} itself at any bidegree.

This release replaces that foundation.  The observation is that
\eqref{eq:degree-difference} is the shadow of a stronger, simpler, and
integral statement about multiplication alone.  Order the coefficient
coordinates as $a_0,\dots,a_r,b_0,\dots,b_s$ and let
$M$ be the $(r+s+1)\times(r+s+2)$ Jacobian matrix of
$m(A,B)=AB$.  Let
\[
 \kappa=(a_0,\dots,a_r,\,-b_0,\dots,-b_s)
\]
be the infinitesimal relative-scaling direction
$(A,B)\mapsto(\lambda A,\lambda^{-1}B)$, which spans $\ker dm$ on the
coprime locus.

\begin{theorem}[Bordered Jacobian identity]\label{thm:bordered}
Let $r,s\ge1$.  For every row vector
$v=(v_0,\dots,v_{r+s+1})$,
\[
 \boxed{\ \det\begin{pmatrix}M\\ v\end{pmatrix}
 =(-1)^{s(r+1)+1}\,\Res(A,B)\,\langle v,\kappa\rangle\ }
\]
as an identity in $\Z[a_0,\dots,b_s,v_0,\dots,v_{r+s+1}]$.
Equivalently, the signed maximal minors of $M$ satisfy
\[
 (-1)^k\det M_{\hat k}
 =(-1)^{r(s+1)}\,\Res(A,B)\,\kappa_k,
 \qquad 0\le k\le r+s+1,
\]
where $M_{\hat k}$ deletes the $k$-th column.  Both statements hold
for all $r,s\ge1$, including $r=s$.
\end{theorem}

The identity should be read as follows: \emph{the top exterior power
of the multiplication differential is the resultant times the
contraction of the source volume form along the lost scaling
direction}.  Every Jacobian obtained by adjoining one scalar function
to multiplication is then computed by a single pairing.

\begin{corollary}[Contraction principle]\label{cor:contraction}
Let $g\in\Z[a_0,\dots,b_s]$ and write
$\Phi_g=(m,g):V_r\times V_s\to V_{r+s}\times\A^1$.  Then
\[
 \det D\Phi_g=(-1)^{s(r+1)+1}\,\Res(A,B)\cdot dg(\kappa).
\]
If $g$ is bihomogeneous of weight $(p,q)$, meaning
$g(\lambda A,\mu B)=\lambda^p\mu^q\,g(A,B)$, then
$dg(\kappa)=(p-q)\,g$ and
\[
 \det D\Phi_g=(-1)^{s(r+1)+1}(p-q)\,g\,\Res .
\]
\end{corollary}

Three specialisations organise the whole subject.

\begin{enumerate}[leftmargin=2.2em]
\item \textbf{The degree-difference identity.}  The resultant is
bihomogeneous of weight $(s,r)$, so $g=\Res$ gives
$\det D\Phi_{r,s}=(-1)^{s(r+1)+1}(s-r)\Res^2
=(-1)^{s(r+1)}(r-s)\Res^2$: identity
\eqref{eq:degree-difference}, now proved over $\Z$
(\cref{cor:integral-dd}).
\item \textbf{The pullback obstruction.}  If $g=h\circ m$ factors
through the product, then $dg(\kappa)=dh(dm(\kappa))=0$, so
$\det D\Phi_g=0$ identically.  No reparametrisation of the product
alone can complete multiplication to a nondegenerate chart; the added
coordinate must break the scaling symmetry
(\cref{cor:pullback}).
\item \textbf{The equal-degree collapse, localised.}
\Cref{thm:bordered} does not degenerate at $r=s$: the minors still
equal $\pm\Res\cdot\kappa_k$.  What vanishes is the Euler weight
$s-r$ of the resultant.  The failure of $\Phi_{r,r}$ recorded in
\cite{DDAS} is thus not a failure of the factorisation geometry; it is
the statement that at equal degrees the resultant becomes invariant
under the scaling torus, so its derivative can no longer see the lost
direction (\cref{rem:equal-degree}).
\end{enumerate}

Together (1)--(3) address, at the infinitesimal level, a question Tao
raised when digesting the counterexample---he wrote that he did
not have ``a completely satisfactory geometric explanation for this
miracle'' \cite{TaoDigestion}.  Among candidate normalising
functions, coefficient functionals pair with $\kappa$ to single
coefficients, which vanish on divisors meeting the coprime locus;
functions of the product pair to zero; the resultant pairs to
$(s-r)\Res$, which is nonvanishing on the \emph{entire} coprime locus
precisely when $r\ne s$.  The resultant is not literally unique in
this: in characteristic zero the normalisers are exactly the
$g=c'\Res^n+h$ with $h$ in the kernel of the scaling derivation, and
only the pullback part of that kernel is absorbed by target
automorphisms (\cref{prop:normalisers}).  The degree difference is
the eigenvalue by which each power of the resultant reproduces itself
under the scaling symmetry.

One limitation must be stated plainly.  In its source context, Tao's
remark concerns the \emph{global} phenomenon: why the repeated-root
slice itself is affine three-space.  The present identity explains
the local \'etaleness and normalisation component of the
construction; it does not explain that global fact, which rests on
the slice geometry of \cite{DDAS}.

\subsection*{What is gained over the July proof}

The proof of \cref{thm:bordered} in \cref{sec:proof} is elementary
linear algebra: interpolation at the roots of $A$ and $B$ plus one
auxiliary point, a gap-Vandermonde identity, the product formula for
the resultant, and Zariski density.  In particular:

\begin{itemize}[leftmargin=2.2em]
\item no irreducibility of the resultant hypersurface is used;
\item no unique factorisation or normality is used;
\item every sign is computed, and each signed constant is
independently machine-verified per bidegree;
\item the identity lives in $\Z[a,b]$, so reduction modulo $p$ is
legitimate: $\Phi_{r,s}$ is \'etale on the coprime locus over any
field of characteristic $p$ if and only if $p\nmid s-r$
(\cref{cor:char-p}), sharpening the characteristic-$p$ remark of
\cite{DDAS} into a theorem.
\end{itemize}

Conceptually, \cref{thm:bordered} is the coordinate form of a
classical principle: the resultant is the determinant of the
degree-$(r+s)$ strand of the Koszul complex of $(A,B)$, a construction
going back to Cayley and developed in modern form by Jouanolou,
Chardin and Demazure
\cite{Cayley1848,Jouanolou1991,Chardin1992,Demazure2012,GKZ}.  The
strand just above the Sylvester strand is exact with one-dimensional
kernel spanned by $\kappa$, and its determinant is again the
resultant.  We have not found the bordered form with its two sign laws
stated in the literature; the search boundary is recorded in
\cref{sec:provenance}.  We make no claim of essential novelty for the
principle itself, and every consequence drawn in
\cref{sec:consequences} either strengthens or re-derives a statement
of \cite{DDAS} on a firmer footing.

\subsection*{Machine verification and formalization}

The accompanying suite \texttt{verify\_\allowbreak bordered\_\allowbreak jacobian.py}
verifies
each identity of this paper per bidegree, exactly over $\Z$ or
$\Q$, in two separately written computational backends---SymPy symbolic
determinants and FLINT integer-polynomial arithmetic with
fraction-free elimination---together with negative controls that
confirm the checks can fail (\cref{sec:verification}).  The full
determinant identity is verified for every bidegree with
$r+s\le8$ and at the spot checks $(4,5)$, $(5,5)$, $(5,6)$; the
bordered/minor identity, including $r=s$, through $r+s\le7$; the
base-point sign matrix for all $r,s\le40$.

The path to a complete machine-checked proof is shorter than one
might expect: Mathlib's \texttt{UniversalFactorizationRing}
development already proves that the Jacobian matrix of the universal
monic factorisation presentation \emph{is} the Sylvester matrix up to
sign---a square truncation of our matrix $M$---and that the universal
coprime factorisation ring is \'etale over any base
\cite{MathlibUFR,MathlibResultant}.  The monic normalisation removes
the scaling direction, so these results are an ingredient rather
than an instance of \cref{thm:bordered}; the degree-difference
content itself remains to be formalized.  \Cref{sec:blueprint} states
the relationship precisely and lists the missing lemmas.

\section{Conventions and elementary lemmas}\label{sec:conventions}

Work over a commutative ring; all identities are in polynomial rings
over $\Z$ and are proved by verification over $\C$ on a dense open
set.  Fix coordinates $X,Y$ and write
$A=\sum_{i=0}^ra_iX^{r-i}Y^i$, $B=\sum_{j=0}^sb_jX^{s-j}Y^j$.
Dehomogenising by $t=Y/X$ gives
$a(t)=\sum a_it^i$ and $b(t)=\sum b_jt^j$.  The product
$C=AB$ has coefficients $c_k=\sum_{i+j=k}a_ib_j$.

\subsection{Sylvester convention}
$\Res(A,B)$ is the determinant of the $(r+s)\times(r+s)$ matrix whose
rows are the coefficient vectors of
$t^ka(t)$ for $0\le k<s$ followed by those of $t^kb(t)$ for
$0\le k<r$, in the basis $1,t,\dots,t^{r+s-1}$.  Then
$\Res(X^r,Y^s)=1$, matching \cite{DDAS}.

\subsection{Jacobian data}
$M$ denotes the $(r+s+1)\times(r+s+2)$ matrix with rows indexed by
$c_0,\dots,c_{r+s}$ and columns by $a_0,\dots,a_r,b_0,\dots,b_s$:
\[
 M[k,i]=\frac{\partial c_k}{\partial a_i}=b_{k-i},
 \qquad
 M[k,r+1+j]=\frac{\partial c_k}{\partial b_j}=a_{k-j},
\]
with out-of-range subscripts read as zero.  Column indices run from
$k=0$ (the $a_0$ column) to $k=r+s+1$ (the $b_s$ column).

\subsection{Vandermonde notation}
For points $x_1,\dots,x_m$ write
$\Vand(x_1,\dots,x_m)=\prod_{p<q}(x_q-x_p)$.  For
$0\le k\le m$ let $W_{\hat k}(x_1,\dots,x_m)$ be the $m\times m$
matrix whose rows are
$(x_p^{c})_{c\in\{0,\dots,m\}\setminus\{k\}}$.

\begin{lemma}[Gap-Vandermonde]\label{lem:gap-vandermonde}
$\det W_{\hat k}(x_1,\dots,x_m)=e_{m-k}(x_1,\dots,x_m)\,
\Vand(x_1,\dots,x_m)$, where $e_j$ is the elementary symmetric
polynomial.
\end{lemma}

\begin{proof}
Let $T$ be an auxiliary variable and consider the full
$(m+1)\times(m+1)$ Vandermonde matrix on the points
$(x_1,\dots,x_m,T)$ with powers $0,\dots,m$.  Its determinant is
\[
 \Vand(x)\prod_p(T-x_p)
 =\Vand(x)\sum_k(-1)^{m-k}e_{m-k}(x)\,T^k .
\]
Expanding instead along the $T$-row gives
$\sum_k(-1)^{m+k}T^k\det W_{\hat k}(x)$.  Comparing the coefficients
of $T^k$ gives the claim.
\end{proof}

\begin{lemma}[Product formula in this convention]\label{lem:product}
Suppose $a(t)=a_r\prod_{i=1}^r(t-\alpha_i)$ and
$b(t)=b_s\prod_{j=1}^s(t-\beta_j)$.  Then
\[
 \Res(A,B)=a_r^s\,b_s^r\prod_{i,j}(\beta_j-\alpha_i)
 =b_s^r\prod_j a(\beta_j)
 =(-1)^{rs}a_r^s\prod_i b(\alpha_i).
\]
\end{lemma}

\begin{proof}
The transpose of the Sylvester matrix represents
$\Psi:(u,w)\mapsto ua+wb$ on
$V_{s-1}\times V_{r-1}\to V_{r+s-1}$ in monomial bases.  Assume the
$\alpha_i,\beta_j$ are pairwise distinct and $a_rb_s\ne0$; the general
case follows by density.  Postcompose with evaluation $E$ at the
$r+s$ points $(\alpha_1,\dots,\alpha_r,\beta_1,\dots,\beta_s)$;
$\det E=\Vand(\alpha,\beta)$.  In the evaluated basis,
$(E\Psi)(u,w)$ has value $w(\alpha_i)b(\alpha_i)$ at $\alpha_i$ and
$u(\beta_j)a(\beta_j)$ at $\beta_j$, so $E\Psi$ is block
anti-diagonal: an $r\times r$ block
$\bigl(b(\alpha_i)\alpha_i^{c}\bigr)_{c<r}$ against the $w$-columns and
an $s\times s$ block
$\bigl(a(\beta_j)\beta_j^{c}\bigr)_{c<s}$ against the $u$-columns.
Hence
\[
 \det(E\Psi)=(-1)^{rs}\,
 \Bigl(\prod_ib(\alpha_i)\Bigr)\Vand(\alpha)\,
 \Bigl(\prod_ja(\beta_j)\Bigr)\Vand(\beta).
\]
Dividing by
$\det E=\Vand(\alpha)\Vand(\beta)\prod_{i,j}(\beta_j-\alpha_i)$
and substituting
$\prod_ib(\alpha_i)=b_s^r\prod_{i,j}(\alpha_i-\beta_j)$
yields the displayed formulas.  At $(A,B)=(X^r,Y^s)$ the matrix is the
identity, confirming the normalisation.
\end{proof}

\begin{lemma}[Euler identity]\label{lem:euler}
$\Res(\lambda A,\mu B)=\lambda^s\mu^r\Res(A,B)$, and consequently
$d\Res(\kappa)=(s-r)\Res$.  More generally, if $g$ is bihomogeneous of
weight $(p,q)$ then $dg(\kappa)=(p-q)g$.
\end{lemma}

\begin{proof}
Scaling $a$ multiplies the $s$ upper Sylvester rows by $\lambda$;
scaling $b$ multiplies the $r$ lower rows by $\mu$.  For the second
claim, write
$\delta=\sum_ia_i\partial_{a_i}-\sum_jb_j\partial_{b_j}$ for the
\emph{scaling derivation}, so that $dg(\kappa)=\delta g$ for every
$g$ by the definition of $\kappa$.  On a monomial of bidegree
$(p,q)$ in the $a$- and $b$-variables, $\delta$ acts as
multiplication by the integer $p-q$; hence
$\delta g=(p-q)\,g$ for bihomogeneous $g$.  This is a formal
identity of polynomials, valid over any commutative ring.
\end{proof}

\begin{remark}[Worked example: $(r,s)=(1,1)$]\label{ex:11}
For $A=a_0X+a_1Y$, $B=b_0X+b_1Y$: the product has coefficients
$c_0=a_0b_0$, $c_1=a_0b_1+a_1b_0$, $c_2=a_1b_1$;
$R=a_0b_1-a_1b_0$; $\kappa=(a_0,a_1,-b_0,-b_1)$; and, on the
column order $(a_0,a_1,b_0,b_1)$,
\[
 M=\begin{pmatrix}
 b_0&0&a_0&0\\
 b_1&b_0&a_1&a_0\\
 0&b_1&0&a_1
 \end{pmatrix},
 \qquad \sigma=(-1)^{r(s+1)}=+1 .
\]
Deleting the $a_0$ column: $\det M_{\hat0}=a_0R$, so
$(-1)^0\det M_{\hat0}=\Res\cdot\kappa_0$.  Deleting the
$b_1$ column: $\det M_{\hat3}=b_1R$, so
$(-1)^3\det M_{\hat3}=-b_1R=\Res\cdot\kappa_3$.  Bordering by any
$v$:
$\det\bigl(\begin{smallmatrix}M\\v\end{smallmatrix}\bigr)
=-R\,(v_0a_0+v_1a_1-v_2b_0-v_3b_1)$, the case
$(-1)^{s(r+1)+1}=-1$ of \cref{thm:bordered}.  Every sign convention
of the paper can be audited on this display.
\end{remark}

\section{Proof of the bordered Jacobian identity}\label{sec:proof}

By multilinearity of the determinant in the border row $v$, the
bordered form of \cref{thm:bordered} is equivalent to the minor form:
expanding along the last row,
\[
 \det\begin{pmatrix}M\\e_k\end{pmatrix}
 =(-1)^{(r+s+1)+k}\det M_{\hat k},
\]
and $(-1)^{r+s+1}\cdot(-1)^{r(s+1)}=(-1)^{s(r+1)+1}$, so the two
displayed sign laws match.  It therefore suffices to prove, for each
column index $k$,
\begin{equation}\label{eq:minor}
 (-1)^k\det M_{\hat k}=(-1)^{r(s+1)}\,\Res(A,B)\,\kappa_k .
\end{equation}
Both sides are polynomials in $\Z[a,b]$, so it suffices to verify
\eqref{eq:minor} on the dense open set
$U\subset\C^{r+s+2}$ where $a_rb_s\ne0$ and the combined list of
roots $\alpha_1,\dots,\alpha_r$ of $a$ and $\beta_1,\dots,\beta_s$ of
$b$ is pairwise distinct, within and across the two lists; $U$ is a
non-empty Zariski-open set.  For the computation, an auxiliary point
$\tau\in\C$ is chosen pointwise, distinct from all roots (such
$\tau$ exist because $\C$ is infinite); $\tau$ enters only through
the evaluation basis, every $\tau$-dependent factor cancels, and the
final expressions are independent of $\tau$.

\subsection{Evaluation}
$M_{\hat k}$ is square of size $r+s+1$.  Change basis on the target
(row space) by evaluating degree-$(r+s)$ polynomials at the
$r+s+1$ points $(\alpha_1,\dots,\alpha_r,\beta_1,\dots,\beta_s,\tau)$;
the change-of-basis determinant is
$\Vand(\alpha,\beta,\tau)$.  Writing $V_E$ for the evaluation matrix,
the rows of $V_EM_{\hat k}$ are, by the computation
$\sum_kx^k\,\partial c_k/\partial a_i=x^ib(x)$ and
$\sum_kx^k\,\partial c_k/\partial b_j=x^ja(x)$:
\begin{align*}
 \text{row }\alpha_i&:\quad
 b(\alpha_i)\cdot\bigl(\alpha_i^{c}\ \big|\ 0\bigr),
 &&c\text{ ranging over the surviving $a$-powers},\\
 \text{row }\beta_j&:\quad
 a(\beta_j)\cdot\bigl(0\ \big|\ \beta_j^{c}\bigr),
 &&c\text{ ranging over the surviving $b$-powers},\\
 \text{row }\tau&:\quad
 \bigl(b(\tau)\,\tau^{c}\ \big|\ a(\tau)\,\tau^{c}\bigr).
\end{align*}

\subsection{Splitting the auxiliary row}
Split the $\tau$-row as
$(b(\tau)\tau^{c}\,|\,0)+(0\,|\,a(\tau)\tau^{c})$ and expand the
determinant additively.  In each of the two resulting matrices, count
rows against columns: if the deleted column $k$ lies in the $b$-block,
the $a$-block retains $r+1$ columns, which only the $r$ rows
$\alpha_i$ and possibly the $\tau$-row can match; the term placing
$\tau$ in the $b$-block therefore has $r+1$ columns to fill with $r$
rows and vanishes.  Symmetrically, if $k$ lies in the $a$-block, the
term placing $\tau$ in the $a$-block vanishes.  Exactly one term
escapes this block-rank vanishing in each case; at special points the
surviving term may of course itself vanish, exactly when the
corresponding coefficient does.

\subsection{The \texorpdfstring{$b$}{b}-block case}
Let $k=r+1+j$ with $0\le j\le s$, so the $b$-column of $b_j$ is
deleted.  The surviving term is block-diagonal after moving the
$\tau$-row above the $\beta$-rows, a transposition past $s$ rows
contributing $(-1)^s$:
\[
 \det\bigl(V_EM_{\hat k}\bigr)
 =(-1)^s\,
 \Bigl[\prod_ib(\alpha_i)\Bigr]b(\tau)\,\Vand(\alpha,\tau)\cdot
 \Bigl[\prod_ja(\beta_j)\Bigr]\det W_{\hat\jmath}(\beta),
\]
where the $a$-block is the full Vandermonde on
$(\alpha_1,\dots,\alpha_r,\tau)$ with powers $0,\dots,r$ and the
$b$-block is the gap matrix of \cref{lem:gap-vandermonde}.  Divide by
$\Vand(\alpha,\beta,\tau)
=\Vand(\alpha)\Vand(\beta)\prod_{i,j}(\beta_j-\alpha_i)
\prod_i(\tau-\alpha_i)\prod_j(\tau-\beta_j)$
and cancel in three moves:
$\Vand(\alpha,\tau)=\Vand(\alpha)\prod_i(\tau-\alpha_i)$ absorbs the
$\alpha$--$\tau$ factors; $b(\tau)=b_s\prod_j(\tau-\beta_j)$ absorbs
the $\beta$--$\tau$ factors, leaving $b_s$; and
$\prod_ib(\alpha_i)=(-1)^{rs}b_s^r\prod_{i,j}(\beta_j-\alpha_i)$
absorbs the mixed factors, leaving $(-1)^{rs}b_s^r$.  With
$\det W_{\hat\jmath}(\beta)=e_{s-j}(\beta)\Vand(\beta)$
(\cref{lem:gap-vandermonde}) this collects to
\[
 \det M_{\hat k}
 =(-1)^{s(r+1)}\,b_s\,e_{s-j}(\beta)\cdot
 b_s^{r}\prod_ja(\beta_j)
 =(-1)^{s(r+1)}\,b_s\,e_{s-j}(\beta)\,\Res,
\]
the last step by the second form of \cref{lem:product}.
Since $b_j=(-1)^{s-j}b_s\,e_{s-j}(\beta)$, this reads
$\det M_{\hat k}=(-1)^{s(r+1)}(-1)^{s-j}\,b_j\,\Res$, and hence
\[
 (-1)^k\det M_{\hat k}
 =(-1)^{r+1+j}(-1)^{s(r+1)+s-j}\,b_j\Res
 =(-1)^{r(s+1)}\,(-b_j)\,\Res
 =(-1)^{r(s+1)}\,\Res\,\kappa_k .
\]

\subsection{The \texorpdfstring{$a$}{a}-block case}
Let $k=i$ with $0\le i\le r$.  Now the term with the $\tau$-row in
the $b$-block survives, and the row order
$(\alpha\text{'s},\beta\text{'s},\tau)$ is already block-ordered, so
no transposition sign appears:
\[
 \det\bigl(V_EM_{\hat k}\bigr)
 =\Bigl[\prod_ib(\alpha_i)\Bigr]\det W_{\hat\imath}(\alpha)\cdot
 \Bigl[\prod_ja(\beta_j)\Bigr]a(\tau)\,\Vand(\beta,\tau).
\]
The same cancellations ($a(\tau)=a_r\prod_i(\tau-\alpha_i)$ now
supplies the $\prod_i(\tau-\alpha_i)$) give
\[
 \det M_{\hat k}
 =(-1)^{rs}\,a_r\,e_{r-i}(\alpha)\,\Res
 =(-1)^{rs}(-1)^{r-i}\,a_i\,\Res,
\]
whence
$(-1)^k\det M_{\hat k}=(-1)^{rs+r}\,a_i\Res
=(-1)^{r(s+1)}\Res\,\kappa_k$, as $\kappa_i=+a_i$.

\subsection{Conclusion}
Identity \eqref{eq:minor} holds on the dense set $U$, hence in
$\Z[a,b]$; nothing in the computation distinguished $r=s$.  This
proves \cref{thm:bordered}.  \qed

\begin{remark}\label{rem:koszul}
Let $S=\Z[X,Y]$ with its standard grading, and order the Koszul
generators as $(B,A)$, so that the first differential in degree
$r+s$ is $M:(\dot A,\dot B)\mapsto\dot AB+\dot BA$ and the second
is $\kappa$, matching the sign convention
$\kappa=(A,-B)$.  The matrix
$\bigl(\begin{smallmatrix}M\\v\end{smallmatrix}\bigr)$ is then the
degree-$(r+s)$ strand of the Koszul complex of $(A,B)$, based and
bordered: the strand
$0\to S_0\xrightarrow{\ \kappa\ }S_r\oplus S_s
\xrightarrow{\ M\ }S_{r+s}\to0$
is exact on the coprime locus because the quotient
$S/(A,B)$ by a coprime pair vanishes in degrees
above $r+s-2$, and \cref{thm:bordered} says that the determinant of
this based exact complex---not of any single differential---is the
resultant, as it is for every strand of degree at least
$r+s-1$: the Cayley determinant
\cite{Cayley1848,Jouanolou1991,Chardin1992,Demazure2012}.  The proof
above is nothing but an explicit basing of that strand by
interpolation.  The bordering technique itself also has a classical
antecedent: for $n-1$ generic \emph{linear} forms, Jouanolou proves
that the signed maximal minors $\Delta_j$ of the coefficient matrix
are the partial derivatives of the bordered determinant
$\Res(l_1,\dots,l_{n-1},\lambda)$ and that
$\Res(l_1,\dots,l_{n-1},f)=f(\Delta)$
\cite[5.4.1--5.4.4]{Jouanolou1991}---bordered evaluation at the
Cramer point, with no resultant factor.  Chardin records the scalar
shadow of the minor form: the resultant divides every maximal minor
of the degree-$\nu$ first Koszul differential for
$\nu>\sum_i(d_i-1)$, and is their gcd in the generic case
\cite[\S IV, Remark~1]{Chardin1992}.  Jouanolou's formulary
develops a bordered-determinant calculus for square elimination
matrices---bordering a Macaulay-type matrix by a row $u$ and a
column $v$ computes the apolarity functional,
$\det\Theta^*_\nu(f;u,v)=c_\nu(f)\,\omega(\overline{uv})$---and
gradient-of-resultant identities modulo $R$ for the square system,
going back to \cite[4.6.3, 6.4.2--6.4.3]{Jouanolou1991}: the
gradient of $R$ with respect to distinguished coefficients
reproduces the common zero mod $R$
\cite[3.10.15, 3.10.25--3.10.31, 3.11.19.23--3.11.19.31]{JouanolouFormulaire}.
The same corpus contains the square-system counterpart of a
Jacobian--resultant identity: for $n$ forms in $n$ variables the
Jacobian determinant with respect to the \emph{variables} satisfies
$\operatorname{Jac}(f_1,\dots,f_n)\equiv
d_1\cdots d_n\,\det(f_{ij})$ modulo $(f_1,\dots,f_n)$, and the
residue functional sends this twisted jacobian to the resultant
\cite[3.8.1.9, 3.8.2.9]{JouanolouAspects}; the linear case of the
signed-minor calculus, with $\Delta_i=\partial\Delta/\partial
U_{1i}$ and a Rees-algebra structure modulo $R$, appears at
\cite[3.5.5]{JouanolouAspects}.  \Cref{thm:bordered} differs from
all of these in differentiating with respect to the
\emph{coefficients} of a \emph{factorisation} map, and in holding
exactly over $\Z$ rather than modulo the ideal or through the
residue pairing; there the natural constant is the \emph{product}
of the degrees, here it is their \emph{difference}.  Against
Chardin's divisibility statement in particular, it is the upgrade
from ``$\Res$ divides every maximal minor, and is their gcd
generically'' to an exact vector identity: each signed minor is
identified as $\Res$ times the corresponding coordinate of the
kernel vector, with both sign laws computed, and the Euler pairing
then produces the degree difference.  We prefer the bordered formulation because it
computes, in one stroke, every Jacobian in which the missing
coordinate might be sought.
\end{remark}

\section{Consequences}\label{sec:consequences}

\begin{corollary}[Integral degree-difference identity]
\label{cor:integral-dd}
In $\Z[a_0,\dots,b_s]$,
\[
 \det D\Phi_{r,s}=(-1)^{s(r+1)}(r-s)\Res(A,B)^2 .
\]
\end{corollary}

\begin{proof}
Border $M$ with $v=d\Res$ and contract:
$\det D\Phi_{r,s}=(-1)^{s(r+1)+1}\Res\cdot d\Res(\kappa)
=(-1)^{s(r+1)+1}(s-r)\Res^2$ by \cref{lem:euler}.
\end{proof}

\begin{corollary}[{\'Etaleness over fields and over $\Z[1/(s-r)]$}]\label{cor:char-p}
Over any field $F$, the map $\Phi_{r,s}$ is \'etale at every point
of the open subscheme $\{\Res\ne0\}$ (all scheme points) if and
only if
$(s-r)\cdot1_F\ne0$ in $F$: in characteristic zero, exactly when
$r\ne s$; in characteristic $p>0$, exactly when
$p\nmid s-r$.  For $r\ne s$, $\Phi_{r,s}$ is \'etale on the coprime
locus over $\Z[1/(s-r)]$.
\end{corollary}

\begin{proof}
The determinant is a unit times $(r-s)\Res^2$ at such points.
Conversely, if $(s-r)\cdot1_F=0$ the determinant vanishes identically
over $F$.
\end{proof}

\begin{corollary}[Pullback obstruction]\label{cor:pullback}
If $g=h\circ m$ for any polynomial $h$ on $V_{r+s}$, then
$\det D(m,g)=0$ identically.  In particular no coefficient of the
product, and no discriminant-type function of the product, can serve
as the added coordinate: these fail identically, everywhere.  By
contrast, a source-coefficient functional such as $a_0$ fails only
on a divisor (\cref{rem:equal-degree}).
\end{corollary}

\begin{remark}[Equal degrees]\label{rem:equal-degree}
At $r=s$ the minor identity survives with sign
$(-1)^{r(s+1)}=+1$: the \emph{signed} maximal minors
$(-1)^k\det M_{\hat k}$ equal
$\Res\cdot\kappa_k$ on the nose, and $M$ has full rank $2r+1$
wherever $\Res\ne0$.  The equal-degree collapse of
\eqref{eq:degree-difference} is therefore located entirely in
\cref{lem:euler}: the resultant has weight $(s,r)=(r,r)$, is
invariant under the scaling torus, and its Euler pairing with
$\kappa$ vanishes.  More generally, by \cref{cor:contraction} no
bihomogeneous $g$ of balanced weight can complete the chart at
$r=s$, while an unbalanced $g$ such as a single coefficient
$a_0$ yields $\det D(m,a_0)=\pm a_0\Res$---nondegenerate off a
divisor, but never on the whole coprime locus, since every
coefficient vanishes somewhere on it.  The scaling direction at equal
degrees is invisible to every function natural for the torus action;
this is the precise infinitesimal content of the statement in
\cite{DDAS} that $\Phi_{r,r}$ is everywhere degenerate.
\end{remark}

\begin{proposition}[Classification of normalisers]
\label{prop:normalisers}
Let $F$ be a field, $S=F[a_0,\dots,a_r,b_0,\dots,b_s]$,
$R=\Res(A,B)\in S$, and let
$\delta=\sum_ia_i\partial_{a_i}-\sum_jb_j\partial_{b_j}$ be the
scaling derivation, so that
$\det D(m,g)=(-1)^{s(r+1)+1}R\,\delta g$ for every $g\in S$ by
\cref{thm:bordered}.  Assume $r\ne s$.
\begin{enumerate}[leftmargin=2.2em,label=\textup{(\roman*)}]
\item $(m,g)$ is \'etale at every point of the open subscheme
$D(R)=\{\Res\ne0\}$ if and only if $\delta g$ is a unit of the
localisation $S[R^{-1}]$, i.e.\ $\delta g=cR^n$ with
$c\in F^\times$ and $n\ge1$.
\item If $\operatorname{char}F=0$, the solutions are exactly
\[
 g=\frac{c}{n(s-r)}\,R^n+h,
 \qquad c\in F^\times,\ n\ge1,\ h\in\ker\delta,
\]
and among bihomogeneous $g$ exactly the constant multiples
$c'R^n$ with $c'\in F^\times$, $n\ge1$.
\item If $\operatorname{char}F=p>0$, the same description holds for
those $n$ with $p\nmid n(s-r)$; the choice $g=cR^n$ fails whenever
$p\mid n(s-r)$, and $\ker\delta$ is spanned by the monomials whose
$\delta$-weight is divisible by $p$, which in general strictly
contains the invariants of the scaling torus.
\end{enumerate}
\end{proposition}

\begin{proof}
By \cref{cor:contraction}, $(m,g)$ is \'etale throughout
$D(R)$ if and only if $R\,\delta g$, equivalently $\delta g$, is a
unit of $\Gamma(D(R),\OO)=S[R^{-1}]$.  Since $S$ is a unique
factorisation domain and $R$ is irreducible over every field
(geometric irreducibility over $\Z$: \cite[2.3(iii)]{Jouanolou1991}),
the units of $S[R^{-1}]$ are $cR^{\Z}$ with $c\in F^\times$; as
$\delta g\in S$, this gives $\delta g=cR^n$ with $n\ge0$.  The
derivation $\delta$ preserves the weight decomposition
$S=\bigoplus_wS_w$ and acts on $S_w$ as multiplication by
$w\cdot1_F$; in particular the weight-zero component of
$\delta g$ vanishes, which excludes $n=0$ and proves (i).  Writing
$g=\sum_wg_w$, the equation $\delta g=cR^n$ reads
$w\,g_w=cR^n\,[w=n(s-r)]$; when $n(s-r)$ is invertible in
$F$ this determines $g_{n(s-r)}=cR^n/(n(s-r))$, leaves
$g_0\in\ker\delta$ free, and forces every other component to
vanish---and when $p\mid n(s-r)$ there is no solution on that
component, while $\delta(cR^n)=n(s-r)cR^n=0$.  A bihomogeneous
$g$ is concentrated in a single weight, which must then be
$n(s-r)\ne0$.  The description of $\ker\delta$ in characteristic
$p$ is immediate from the action on monomials.
\end{proof}

\begin{remark}[The kernel is larger than the pullbacks]
\label{rem:kernel}
Only part of the ambiguity $h\in\ker\delta$ is geometrically
trivial: for a pullback $h=h'\circ m$ the maps $(m,g+h)$ and
$(m,g)$ differ by the target automorphism
$(C,z)\mapsto(C,z+h'(C))$, but $\ker\delta$ is strictly larger than
the pullbacks.  At $r=s=1$, the element $h=a_0b_1-a_1b_0$ satisfies
$\delta h=0$ and $h^2=c_1^2-4c_0c_2$ is a pullback, while $h$ itself
is not: it is antisymmetric under the swap
$(A,B)\mapsto(B,A)$, which fixes every pullback.  For such
$h$, $(m,g+h)$ has the same Jacobian determinant as $(m,g)$ without
being obtained from it by a target automorphism.  The same integer
$s-r$ appears in the \'etaleness criterion, the
$\mu_{\abs{s-r}}$-torsor, the class group
$\Z/\abs{s-r}\Z$ and the covering degree
$\abs{s-r}\binom{r+s}{r}$ of \cite{DDAS} because each of those
statements consumes the determinant identity or the Euler weight as
its local input; the global statements themselves are proved by the
global arguments of \cite{DDAS}, not by the pairing alone.
\end{remark}

\begin{remark}[Effect on the July release]\label{rem:effect}
All downstream results of \cite{DDAS} that consumed identity
\eqref{eq:degree-difference}---\'etaleness of the marked-root map,
the torsor structure, the classification of linear--quadratic slices,
and the Euler-characteristic exclusion of higher-degree untrimmed
slices---now draw their \emph{determinant input} from
\cref{thm:bordered,lem:euler} alone; their global arguments are
unchanged and remain as in \cite{DDAS}.  The dependency audit,
statement by statement (each July item was reread against its text
for this release):

\begin{center}
\small
\begin{tabular}{@{}lll@{}}
\toprule
July statement & determinant input consumed & status \\
\midrule
Thm.~2.2 (identity) & --- & superseded by \cref{cor:integral-dd} \\
\'etaleness on $\{\Res\ne0\}$ & the identity & re-proved over $\Z$ (\cref{cor:char-p}) \\
Prop.~2.4 ($\mu_{\abs{s-r}}$-torsor) & none (Kummer covers) & reread; independent \\
Prop.~2.5 (generic degree) & torsor count & reread; unchanged \\
Prop.~2.7, Cor.~2.8 (units, $\Cl$, $\chi_c$) & torsor + localisation & reread; unchanged \\
Prop.~3.3 (marked-root \'etale) & base change of identity & input replaced \\
\S\S4--6 (classification, tangent slice) & \'etaleness + slice geometry & reread; input via rows above \\
\bottomrule
\end{tabular}
\end{center}

The two inputs the July
proof of the identity needed (irreducibility of the resultant
hypersurface; the base-point inversion count) are no longer
load-bearing \emph{for the identity}: irreducibility remains an input
to the unit classification used in \cref{prop:normalisers} and in
the global theory of \cite{DDAS}, and the base-point computation
survives as a machine-checked consistency anchor.
\end{remark}

\section{Machine verification}\label{sec:verification}

The suite \texttt{verify\_bordered\_jacobian.py} accompanying this
release verifies, exactly and with no floating point or sampling,
the following checks.  The backend of each check is stated.  The two
backends---SymPy symbolic and FLINT exact integer-polynomial with
fraction-free elimination---are separately written construction and
check paths, but they live in one suite, written by the same author
from the same conventions: they guard against implementation error,
not against a shared specification error, and their coefficientwise
cross-comparison applies to the central determinant identity, not to
every item.  Receipts record package versions, platform and
per-check timings; the degree bounds are chosen audit boundaries,
not resource limits, except that $(6,7)$ and beyond were measured
impractical (the $(5,6)$ spot check runs in about six minutes).

\begin{enumerate}[label=(\Roman*),leftmargin=2.6em]
\item[(0)] \cref{lem:product} with symbolic roots, for
$(r,s)\in\{(1,1),(1,2),(2,2),(2,3),(3,3)\}$ [SymPy];
\item[(0$'$)] the gap-Vandermonde identity of
\cref{lem:gap-vandermonde} for $2\le m\le5$, all $k$ [SymPy];
\item \cref{cor:integral-dd} for every bidegree with
$r+s\le8$ and for $(4,5),(5,5),(5,6)$ [FLINT], and independently for
$r+s\le5$ [SymPy], with coefficient-identical agreement of the two
backends for $r+s\le5$;
\item the minor identity \eqref{eq:minor} with its sign law, for all
bidegrees with $r+s\le7$ \emph{including} $r=s$ [FLINT];
\item \cref{cor:contraction} at four bidegrees for six weight cases,
including three vanishing cases: balanced weight, a linear pullback,
and a nonlinear pullback $g=c_0^2+c_1$ [FLINT];
\item the base-point sign matrix for all $1\le r,s\le40$
($1600$ cases) [FLINT integer matrices];
\item the forward direction of \cref{cor:char-p} via divisibility of
every coefficient of $\det D\Phi_{r,s}$ by $p$ for five pairs with
$p\mid r-s$ [FLINT]; the converse direction is immediate from the
identity and is not separately tested;
\item negative controls: a perturbed Jacobian entry, a wrong sign, a
wrong resultant power, a sign-flipped kernel vector, and transposed
bidegrees are each detected as failures [FLINT].
\end{enumerate}

The density step and the all-degree quantifier of
\cref{thm:bordered} are not machine-checkable by instance testing and
are not claimed to be.

Two boundaries are declared explicitly.  First, the per-bidegree
checks certify instances, not the all-degree statement; the
all-degree statement rests on the proof in \cref{sec:proof}.  Second,
the ledger of the July release verified the base-point sign matrix
and the cubic-specific algebra but not identity
\eqref{eq:degree-difference} itself; the present suite closes that
gap and supersedes it.  This claim about the July ledger is auditable
against the archived release \cite{DDAS}: its script
(\texttt{verify\_\allowbreak degree\_\allowbreak difference\_\allowbreak
affine\_\allowbreak slices\_\allowbreak final.py}, SHA-256
\texttt{21dbdc38}\dots\footnote{Full digest:
\texttt{21dbdc3835302a40ed636e2090f42040d0\allowbreak
63a669fb1e37ede62d30be8f07834d}.})
enumerates six checks in its receipt, none of which computes
$\det D\Phi_{r,s}$ symbolically.

\section{Formalization blueprint}\label{sec:blueprint}

Mathlib (master commit \texttt{1f0fbd1a}, 6 August 2026) contains:

\begin{center}
\small
\begin{tabular}{@{}p{2.35in}p{3.35in}@{}}
\toprule
This paper & Mathlib analogue \\
\midrule
Sylvester matrix, $\Res$ &
\texttt{Polynomial.sylvester}, \texttt{Polynomial.resultant}
\cite{MathlibResultant} \\
$\Res$ unit $\Leftrightarrow$ coprime &
\texttt{isUnit\_resultant\_iff\_isCoprime} \\
\cref{lem:product} &
\texttt{resultant\_eq\_prod\_eval} (monic normalisation) \\
Jacobian of monic multiplication is the Sylvester matrix (a square
truncation of $M$; see below) &
\texttt{universalFactorizationMapPresentation\_}\newline
\texttt{jacobiMatrix}
\cite{MathlibUFR} \\
\'Etaleness of monic coprime factorisation (related background; see
below) &
\texttt{UniversalCoprimeFactorizationRing} is \'etale
\cite{MathlibUFR} \\
\bottomrule
\end{tabular}
\end{center}

The relationship between the last two rows and the present theorems
must be stated with care, and is weaker than the word ``chart'' would
suggest.  The universal monic factorisation map fixes
$a_r=b_s=1$ and drops the (automatically monic) top coefficient of
the product; its Jacobian is the square matrix obtained from
$M$ by deleting the $a_r$ and $b_s$ columns and the $c_{r+s}$ row.
Mathlib proves this square matrix is the Sylvester matrix, with
determinant $\pm\Res$, and that the coprime factorisation ring is
\'etale over any base---\emph{with no degree-difference factor}.
Indeed the monic multiplication map is \'etale on its coprime locus
even at $r=s$ (at $(r,s)=(1,1)$ its Jacobian determinant is
$b_0-a_0=-\Res$), because fixing the leading coefficients removes the
scaling direction that \cref{thm:bordered} and
\cref{cor:char-p} are about.  The Mathlib results are therefore an
ingredient---the machine-checked determinant of the truncated
matrix---rather than a monic instance of the present statements.
What remains to formalize for the full \cref{thm:bordered}:
(i) bihomogeneity of the resultant (row-scaling of
\texttt{Polynomial.sylvester}); (ii) the Euler identity
\cref{lem:euler}; (iii) the bordered bookkeeping relating the
non-monic $(r+s+2)$-column Jacobian, with its resultant row, to the
truncated Sylvester block.  All three are identities in polynomial
rings over $\Z$, requiring no schemes and no analysis; (iii) is
finite linear algebra over $\Z[a,b]$, and it---not the Mathlib
material---carries the degree-difference content.  The Keller
disproof itself is already fully formalized: pull request \#4474 to
the Formal Conjectures repository, merged 26 July 2026 at commit
\texttt{393aa9a} \cite{LezeauPR}.  A formalized
\cref{thm:bordered} would therefore not fill a gap in the formal
disproof; it would add a machine-checked \emph{factorisation-space
explanation} of the determinant that the disproof verifies
directly.

\section{Provenance, priority and scope}\label{sec:provenance}

The Keller counterexample was announced by Alp\"oge on 20 July 2026,
crediting Fable for the construction
\cite{AlpogeAnnouncement,AnonymousCounterexample}; Gallagher exhibited
an infinite family and Speyer a tangent-sweep geometric account in the
same week, with expositions by Levinson and Tao
\cite{SpeyerThread,LevinsonGeometry,TaoDigestion}.  Gao has since
generalised the tangent-sweep mechanism to all dimensions
$\ge3$ \cite{Gao2026}; that framework is disjoint from the
factorisation-space approach considered here.  The determinant
identity \eqref{eq:degree-difference} was stated and proved for all
bidegrees, over $\C$ and with the two inputs discussed above, in the
July release \cite{DDAS}; the public threads contain the class-group
observations of Skooi and Mondal cited there, but neither the
identity nor its bordered refinement.  The
$(r,s)=(1,2)$ case, with the same sign, appears as an explicit
polynomial identity in Lou's derivation note of 20 July 2026
\cite[Lemma~1]{Lou2026}; the all-degree statement and the
degree-difference factor are the July release's.  Adjacent
determinantal literature on binary forms includes the Jacobian ideal
of the binary discriminant \cite{DAndreaChipalkatti2007}.

The mathematical core of \cref{thm:bordered}---the resultant as the
determinant of any sufficiently high strand (degree at least
$r+s-1$) of the Koszul complex of two binary forms---is classical: Cayley's 1848 construction, developed by
Jouanolou and Chardin, with textbook treatments by Demazure and by
Gelfand--Kapranov--Zelevinsky
\cite{Cayley1848,Jouanolou1991,Chardin1992,Demazure2012,GKZ}.  The
full texts of \cite{Jouanolou1991}, \cite{JouanolouAspects},
\cite{JouanolouFormulaire} and \cite{Chardin1992} were audited for
this release (6 August 2026).  None contains a strand determinant above the
Sylvester degree for two binary forms, the exact identification of
the signed minors as $\Res\cdot\kappa_k$, the sign laws, the
$\Res^2$ Jacobian, or the degree-difference factor.  Their nearest
contents, catalogued in \cref{rem:koszul}, are: the linear-forms
bordered-evaluation principle \cite[5.4.1--5.4.4]{Jouanolou1991};
gradient-of-resultant identities modulo $R$ for the square
elimination system \cite[4.6.3, 6.4.2--6.4.3]{Jouanolou1991},
developed further in \cite[3.10.15, 3.10.31]{JouanolouFormulaire};
the bordered Macaulay-matrix calculus computing the apolarity
functional \cite[3.10.25--3.10.30,
3.11.19.23--3.11.19.31]{JouanolouFormulaire}; determinantal-ideal
computations for two binary forms \emph{below} the Sylvester degree
\cite[3.11.18.18--3.11.18.35]{JouanolouFormulaire}; the
twisted-jacobian identities and the linear-case signed-minor
calculus \cite[3.5.5, 3.8.1.9, 3.8.2.9]{JouanolouAspects}; and
Chardin's divisibility/gcd remark
\cite[\S IV, Remark~1]{Chardin1992}.
Within these audits and a bounded search of the remaining sources,
the public discussion, and current literature databases, we did not
find the bordered coordinate form with its sign laws
$(-1)^{r(s+1)}$, $(-1)^{s(r+1)+1}$, the uniform treatment of
$r=s$, or the contraction-principle derivation of
\eqref{eq:degree-difference}.  We therefore claim as contributions of
this release: a complete elementary integral proof of
\cref{thm:bordered} with all signs computed; the derivation of the
degree-difference identity, the characteristic-$p$ criterion, the
pullback obstruction and the equal-degree localisation from that
single statement; the two-backend verification suite with
negative controls; and the formalization blueprint.  We claim no
priority over the classical strand principle.  The antecedent
questions for all four primary classical sources ---
\cite{Jouanolou1991}, \cite{JouanolouAspects},
\cite{JouanolouFormulaire} and \cite{Chardin1992} --- are settled
by the full-text audits described above.  The differential content
of the corpus (jacobian and Fitting ideals of the resultant
hypersurface \cite[3.5.3--3.5.4]{JouanolouAspects}; modulo-$R$
gradient congruences \cite[3.7.1]{JouanolouAspects}; the twisted
jacobian \cite[3.8]{JouanolouAspects}) is square-system material,
differentiating with respect to the variables or working modulo
$R$ or through the residue pairing; it contains neither the
bordered form nor the exact coefficient-side identity.  One
reading-scope note: within \cite{Jouanolou1991}, \S\S4.6 and 6.4
are known to us through the corpus's own citations
(\cite[3.5.4, 3.7.1.4]{JouanolouAspects},
\cite[3.10.15]{JouanolouFormulaire}) rather than by direct
reading.  The bounded-search qualification now applies only to the
literature beyond this corpus.

\subsection*{Attribution}
This manuscript, its proofs and its verification suite were produced
by AI systems (Anthropic Claude, with the July release produced
jointly with OpenAI Codex) operating under human direction, as part
of the Evidence Press programme.  The status of this release is:
unrefereed candidate, awaiting independent verification.

\begingroup
\footnotesize
\setlength{\parskip}{0pt}
\begin{thebibliography}{99}

\bibitem{AlpogeAnnouncement}
L.~Alp\"oge,
\emph{\href{https://x.com/__alpoge__/status/2079028340955197566}
{Hello there the Jacobian conjecture is false}},
public announcement on X, 20 July 2026.

\bibitem{AnonymousCounterexample}
Anonymous,
\emph{\href{https://www.ulam.ai/research/jacobian.pdf}
{A counterexample to the Jacobian conjecture}},
manuscript dated 20 July 2026.

\bibitem{Cayley1848}
A.~Cayley,
\emph{On the theory of elimination},
Cambridge and Dublin Mathematical Journal \textbf{3} (1848), 116--120.

\bibitem{Chardin1992}
M.~Chardin,
\emph{The resultant via a Koszul complex},
in: Computational Algebraic Geometry (MEGA 1992),
Progr. Math. \textbf{109}, Birkh\"auser, 1993, 29--39.

\bibitem{DDAS}
Anonymous,
\emph{\href{https://doi.org/10.5281/zenodo.21647593}
{The degree-difference principle and affine slices of binary-form
factorisation spaces}},
Evidence Press release, 27 July 2026.
DOI: 10.5281/zenodo.21647593.

\bibitem{DAndreaChipalkatti2007}
C.~D'Andrea and J.~Chipalkatti,
\emph{On the Jacobian ideal of the binary discriminant} (with an
appendix by A.~Abdesselam),
Collect. Math. \textbf{58} (2007), 155--180.

\bibitem{Demazure2012}
M.~Demazure,
\emph{R\'esultant, discriminant},
Enseign. Math. (2) \textbf{58} (2012), 333--373.

\bibitem{Gao2026}
S.~Gao,
\emph{\href{https://arxiv.org/abs/2608.00222}
{Counterexamples to the Jacobian conjecture in dimensions greater
than two}},
arXiv:2608.00222, August 2026.

\bibitem{GKZ}
I.~M. Gelfand, M.~M. Kapranov and A.~V. Zelevinsky,
\emph{Discriminants, Resultants, and Multidimensional Determinants},
Birkh\"auser, 1994.

\bibitem{Jouanolou1991}
J.-P. Jouanolou,
\emph{Le formalisme du r\'esultant},
Adv. Math. \textbf{90} (1991), 117--263.

\bibitem{JouanolouAspects}
J.-P. Jouanolou,
\emph{Aspects invariants de l'\'elimination},
Adv. Math. \textbf{114} (1995), 1--174.

\bibitem{JouanolouFormulaire}
J.-P. Jouanolou,
\emph{Formes d'inertie et r\'esultant: un formulaire},
Adv. Math. \textbf{126} (1997), 119--250.

\bibitem{LevinsonGeometry}
J.~Levinson,
\emph{\href{https://levjake.wordpress.com/2026/07/22/the-jacobian-conjecture-counterexample/}
{The Jacobian conjecture counterexample, via classical projective
geometry}},
22 July 2026.

\bibitem{LezeauPR}
P.~Lezeau,
\emph{\href{https://github.com/google-deepmind/formal-conjectures/pull/4474}
{feat: add Jacobian disproof}},
pull request \#4474, google-deepmind/formal-conjectures, merged
26 July 2026, commit \texttt{393aa9ab2403a2fb}.

\bibitem{Lou2026}
A.~Lou,
\emph{\href{https://aaronlou.com/jacobian_counterexample_derivation.pdf}
{Deriving an explicit polynomial counterexample to the Jacobian
conjecture}},
note dated 20 July 2026 (accessed 6 August 2026).

\bibitem{MathlibResultant}
Mathlib contributors,
\emph{\href{https://leanprover-community.github.io/mathlib4_docs/Mathlib/RingTheory/Polynomial/Resultant/Basic.html}
{Mathlib.RingTheory.Polynomial.Resultant.Basic}},
Mathlib4, accessed 6 August 2026.

\bibitem{MathlibUFR}
Mathlib contributors,
\emph{\href{https://leanprover-community.github.io/mathlib4_docs/Mathlib/RingTheory/Polynomial/UniversalFactorizationRing.html}
{Mathlib.RingTheory.Polynomial.UniversalFactorizationRing}},
Mathlib4, accessed 6 August 2026.

\bibitem{SpeyerThread}
D.~Speyer and contributors,
\emph{\href{https://sbseminar.wordpress.com/2026/07/20/the-new-counterexample-to-the-jacobian-conjecture/}
{The new counterexample to the Jacobian conjecture}},
Secret Blogging Seminar, 20--24 July 2026.

\bibitem{TaoDigestion}
T.~Tao,
\emph{\href{https://terrytao.wordpress.com/2026/07/21/a-digestion-of-the-jacobian-conjecture-counterexample/}
{A digestion of the Jacobian conjecture counterexample}},
21 July 2026, including comment thread.

\end{thebibliography}
\endgroup

\end{document}
